\section{Definição do Problema de Controle}
\label{sec:definicao_problema_controle}

O controle é responsável por exercer as ações de propulsão e da manutenção da atitude da espaçonave visitante, bem como os torques das juntas do manipulador robótico, de modo que:
(i) na aproximação final, a orientação relativa seja zero;
(ii) os limites cinemáticos e dinâmicos sejam respeitados;
(iii) a reação na base seja compensada de forma a atender os limites estabelecidos.  
A técnica adotada neste trabalho é o controle Proporcional--Derivativo (PD) aplicado em diferentes malhas, e quando necessário, é realizada a filtragem do termo derivativo e aplicado saturação.

\subsection{Definição de Estados e Variáveis de Controle}
\label{sec:var_estado_controle}

Para este trabalho, um conjunto de variáveis de estado que descrevem a dinâmica orbital, a atitude e os elos do MR foram estabelecidos, e são:

\begin{equation}
    \label{eq:variaveis_estado}
    \mathbf{x} =
    \{ \vec r,\; \vec v,\; q,\; \vec \omega,\; \theta,\; \dot\theta \},
\end{equation}

onde $\vec r$ e $\vec v$ representam a posição e a velocidade relativas, expressas no referencial LVLH;
$q$ corresponde ao quaternion unitário que descreve a atitude da espaçonave visitante em relação ao referencial inercial;
$\vec \omega$ é o vetor de velocidade angular da base da espaçonave;
$\theta$ indica os ângulos articulares do MR;
e $\dot \theta$ representa as respectivas velocidades articulares.
As variáveis de controle são representadas por forças e torques:

\begin{equation}
    \label{eq:variaveis_controle}
    \mathbf{u} =
    \{ \vec f,\; \vec \tau \},
\end{equation}

onde $\vec f$ são as forças de propulsão que atuam no controle orbital, tanto absoluto quanto relativo, e $\vec \tau$ são os torques aplicados para o controle de atitude da espaçonave e para o acionamento das juntas do manipulador.

O modelo dinâmico considerado integra diferentes subsistemas, como:
(i) as equações de Hill/Clohessy--Wiltshire (HCW) linearizadas para a descrição do movimento relativo,
(ii) as equações de atitude baseadas em formulações de Euler e quaternions, e
(iii) as equações dinâmicas do manipulador robótico acoplado à base da espaçonave.

\subsection{Restrições Operacionais}
\label{sec:restricoes_operacionais}

As restrições operacionais definem tanto os limites mecânicos e dinâmicos do manipulador robótico quanto as condições de aproximação entre as espaçonaves.
No caso do MR, o espaço de trabalho (\emph{workspace}) é o alcance geométrico de seus elos e pela rotação permitida em cada junta.
A ponta do MR (\emph{end-effector} ou $\{E\}$) deve permanecer dentro desse espaço de trabalho, não ultrapassando os limites estabelecidos. 
Durante a fase de aproximação final, a trajetória da espaçonave visitante deve respeitar um corredor de chegada seguro (\emph{keep-out zone} ou KOZ).
Esse corredor é modelado como um cone em torno do eixo longitudinal ($R-bar$) da porta de acoplamento, restringindo a aproximação em relação ao eixo principal e também no plano lateral.
Existem também os limites dinâmicos dos atuadores. Cada junta revoluta do manipulador possui um torque máximo, enquanto a junta prismática está sujeita a uma força limite:

\begin{equation}
\label{eq:limite_torque}
|\tau_j| \leq \tau_{j,\max}, \qquad \text{juntas revolutas},
\end{equation}

\begin{equation}
\label{eq:limite_forca}
|f_5| \leq f_{5,\max}, \qquad \text{junta prismática}.
\end{equation}

Na aproximação de proximidade, adota-se também um limite para a velocidade relativa da espaçonave visitante em relação ao alvo, bem como uma distância mínima de segurança antes da entrada no corredor de acoplamento.
\subsection{Critérios de Desempenho}
\label{sec:criterios_desempenho}

A avaliação do controle adotado nesta dissertação é em função das trajetórias relativas, à atitude e ao MR.
Os critérios considerados são:
(i) o erro de posição relativa nos eixos $(x,y,z)$ do referencial LVLH;  
(ii) o erro de velocidade relativa;  
(iii) o tempo de acomodação;  
(iv) sobressinal máximo;  
(v) erro de atitude;  
e (vi) a suavidade do esforço de controle.
A escolha do controlador PD é baseado na simplicidade de implementação, baixa exigência computacional e robustez em condições operacionais.
Por outro lado, métodos alternativos, como o controle ótimo linear (LQR) e o controle preditivo baseado em modelo (MPC), podem reduzir os custos em termos de consumo de energia e da otimização de trajetória, mas possuem maior complexidade de projeto e maior custo computacional.
Sendo assim, para os objetivos desta pesquisa, o controlador PD atende os requisitos básicos de rastreamento e estabilização.

As condições de sucesso da missão foram definidas como:  
(i) a convergência da atitude relativa para zero;  
(ii) velocidade relativa inferior ao limite operacional de captura;  
(iii) ausência de violações do corredor de aproximação (\emph{keep-out zone}) e 
(iv) ausência de violações do espaço de trabalho do manipulador robótico.  



